% !PRINTED 148
\weberpage{148}
%<<CONT

\weq{1}{\begin{aligned}
\frac{\omega_{1}}{2K}\operatorname{sn}\frac{2Ku}{\omega_{1}}&=\frac{\sigma(u)}{\sigma_{01}(u)},\\
\operatorname{cn}\frac{2Ku}{\omega_{1}}&=\frac{\sigma_{10}(u)}{\sigma_{01}(u)},\\
\operatorname{dn}\frac{2Ku}{\omega_{1}}&=\frac{\sigma_{00}(u)}{\sigma_{01}(u)},
\end{aligned}}

oder auch, indem man von der Homogeneität der $\sigma$-Funktion Gebrauch macht [\S\,37\,(8)]:

\weq{2}{\begin{aligned}
\operatorname{sn}v&=\frac{\sigma(v,\,2K,\,2iK')}{\sigma_{01}(v,\,2K,\,2iK')},\\
\operatorname{cn}v&=\frac{\sigma_{10}(v,\,2K,\,2iK')}{\sigma_{01}(v,\,2K,\,2iK')},\\
\operatorname{dn}v&=\frac{\sigma_{00}(v,\,2K,\,2iK')}{\sigma_{01}(v,\,2K,\,2iK')}.
\end{aligned}}

Die auf der rechten Seite von (2) vorkommenden $\sigma$-Funktionen hängen nur von den beiden Variablen $v$, $\omega$ ab und es kommt zunächst darauf an, die Änderung dieser Funktionen bei Anwendung der Fundamentaltransformationen
\[
(\omega,\ \omega+1), \qquad \left(\omega,\ -\frac{1}{\omega}\right)
\]
zu bestimmen. Wegen
\[
\pi\,\vartheta_{00}^{2}=2K, \qquad 2iK'=2\omega K
\]
erhält man aus \S\,31\,(6) und (11) die entsprechenden Änderungen:
\[
\begin{aligned}
&\omega, && 2K, && 2iK', && \varkappa, && \varkappa',\\
&\omega+1, && 2\varkappa'K, && 2\varkappa'(K+iK'), && \frac{i\varkappa}{\varkappa'}, && \frac{1}{\varkappa'},\\
&-\frac{1}{\omega}, && 2K', && 2iK, && \varkappa', && \varkappa.
\end{aligned}
\]
Setzt man für den Augenblick:
\[
f(v,\omega)=\sigma(v,\,2K,\,2iK'),
\]
so ergibt sich:
\[
\begin{aligned}
f(v,\omega+1)&=\sigma\bigl[v,\,2\varkappa'K,\,2\varkappa'(K+iK')\bigr]\\
&=\sigma(v,\,2\varkappa'K,\,2\varkappa'iK')\ \text{[nach \S\,35\,(7)]}\\
&=\varkappa'\,\sigma\!\left(\frac{v}{\varkappa'},\,2K,\,2iK'\right)\ \text{[nach \S\,37\,(8)]},\\
f\!\left(v,-\frac{1}{\omega}\right)&=\sigma(v,\,2K',\,2iK)\\
&=\sigma(v,\,-2iK,\,2K')\ \text{[nach \S\,35\,(7)]}\\
&=-i\,\sigma(iv,\,2K,\,2iK')\ \text{[nach \S\,37\,(8)]},
\end{aligned}
\]
%CONT>>
